% Local editorial additions to VI/26; existing reader bridge and canonical sources are unchanged.
\par\medskip
\noindent\textbf{Worked examples: finite type, finiteness and Noetherianity.}
For an affine morphism $\operatorname{Spec}B\to\operatorname{Spec}A$,
``finite type'' means $B$ is generated by finitely many elements
\emph{as an $A$-algebra}; ``finite'' means $B$ is finitely generated
\emph{as an $A$-module}.  Every finite ring map is of finite type,
but the reverse implication fails.  For general morphisms, finite type
also requires quasi-compactness.  The examples below make the
various finiteness conditions distinguishable by direct calculations.

\par\medskip
\noindent\textbf{Worked example 1 (finite type is not a finite morphism).}
Let $A=k$ be a field and $B=k[z]$.
The single element $z$ generates $B$ as a $k$-algebra:
$B=k[z]$; therefore
$\operatorname{Spec}k[z]\to\operatorname{Spec}k$ is of finite type.
But $B$ is \emph{not} a finite-dimensional $k$-vector space: the
monomials $1,z,z^2,\ldots$ are linearly independent.
By contrast, $B'=k[z]/(z^3)$ is a finite $k$-module with basis
$1,\bar z,\bar z^2$, so $\operatorname{Spec}B'\to\operatorname{Spec}k$
is a finite morphism.  Both rings are Noetherian, despite the first
map not being finite.  The distinction is algebra-generation versus
module-generation, not a distinction between Noetherian and
non-Noetherian coordinate rings.

\par\medskip
\noindent\textbf{Worked example 2 (localization is finite type, not finite).}
Consider the open immersion
$\operatorname{Spec}\mathbb Z[1/2]\to\operatorname{Spec}\mathbb Z$.
Its coordinate ring has the finite presentation
\[
\mathbb Z[1/2]\cong\mathbb Z[s]/(2s-1),
\]
so the morphism is of finite type (indeed of finite presentation).
It is not finite: if fractions $2^{-n}$ were contained in a
finitely generated $\mathbb Z$-submodule, choose the largest
power $2^N$ needed to clear all denominators of the generators.
That submodule would lie in $2^{-N}\mathbb Z$, but
$2^{-(N+1)}\notin2^{-N}\mathbb Z$, a contradiction.
The corresponding image is the basic open $D(2)$; its complement
is the closed point $(2)$.

\par\medskip
\noindent\textbf{Worked example 3 (Noetherianity via Hilbert's basis theorem).}
Let $A=k[t]$ and
$B=k[t,x,y]/(y^2-x^3-t)$.
First $k[t]$ is Noetherian; consequently $k[t,x,y]$ is Noetherian
by two applications of Hilbert's basis theorem.  Its quotient $B$
is also Noetherian.  Moreover, $B$ is generated over $A$ by
$x$ and $y$, so the structure morphism is of finite type.
At a closed point $(t-a)$ of $\operatorname{Spec}A$, the fiber ring is
\[
B\otimes_A k\cong k[x,y]/(y^2-x^3-a).
\]
This computes a whole family of affine plane curves while keeping
the finite-type and Noetherian assertions separate:
Noetherianity concerns ascending chains of ideals in $B$,
whereas finite type concerns the two generators $x,y$ over $A$.

\par\medskip
\noindent\textbf{Worked example 4 (locally finite type need not be finite type).}
Let $X=\coprod_{n\ge1}\operatorname{Spec}k$, a disjoint union
of countably many $k$-points, and map $X\to\operatorname{Spec}k$.
Each component is an affine open subset mapped isomorphically to
$\operatorname{Spec}k$, so the morphism is \emph{locally of finite type}.
However, the infinitely many components form an open cover of $X$
with no finite subcover.  Thus $X$ is not quasi-compact, and the
morphism is not of finite type.
This counterexample isolates the additional quasi-compactness
requirement; checking finite generation on each affine chart alone
is not sufficient for a global finite-type statement.

\par\medskip
\noindent\textbf{Worked example 5 (the base must be Noetherian).}
Let $A$ be the polynomial algebra over $k$ in countably many
independent variables $x_1,x_2,\ldots$, and put $B=A[u]$.
The ring $B$ is generated over $A$ by the single element $u$,
so $\operatorname{Spec}B\to\operatorname{Spec}A$ is of finite type.
Nevertheless $A$ is not Noetherian, as the strictly increasing
chain
\[
(x_1)\subsetneq(x_1,x_2)\subsetneq(x_1,x_2,x_3)
\subsetneq\cdots
\]
shows.  The same chain, viewed inside $B$, is still strictly
increasing (specialize $u\mapsto0$ to see this), hence $B$ is not
Noetherian either.  The usual theorem that finite-type algebras
over a Noetherian ring are Noetherian needs the Noetherian
hypothesis on the base; finite type alone does not provide it.
